\documentclass[11pt,a4paper]{ctexart}

\usepackage[a4paper,margin=2.35cm]{geometry}
\usepackage{amsmath}
\usepackage{booktabs}
\usepackage{longtable}
\usepackage{array}
\usepackage{enumitem}
\usepackage{xcolor}
\usepackage{listings}
\usepackage{hyperref}
\usepackage{fancyhdr}

\definecolor{codebg}{HTML}{F6F8FA}
\definecolor{codeframe}{HTML}{D0D7DE}
\definecolor{codegreen}{HTML}{116329}
\definecolor{codepurple}{HTML}{8250DF}
\definecolor{codeblue}{HTML}{0550AE}

\lstdefinestyle{rcode}{
  language=R,
  backgroundcolor=\color{codebg},
  frame=single,
  rulecolor=\color{codeframe},
  basicstyle=\ttfamily\small,
  keywordstyle=\color{codeblue}\bfseries,
  commentstyle=\color{codegreen},
  stringstyle=\color{codepurple},
  numbers=left,
  numberstyle=\scriptsize\color{gray},
  numbersep=8pt,
  breaklines=true,
  showstringspaces=false,
  columns=fullflexible,
  keepspaces=true,
  tabsize=2,
  xleftmargin=1em,
  framexleftmargin=0.6em
}

\lstdefinestyle{textcode}{
  backgroundcolor=\color{codebg},
  frame=single,
  rulecolor=\color{codeframe},
  basicstyle=\ttfamily\small,
  numbers=none,
  breaklines=true,
  showstringspaces=false,
  columns=fullflexible,
  keepspaces=true,
  tabsize=2,
  xleftmargin=1em,
  framexleftmargin=0.6em
}

\hypersetup{
  colorlinks=true,
  linkcolor=codeblue,
  urlcolor=codeblue,
  pdftitle={StatProg II Practical 3 -- Debugging},
  pdfauthor={整理版}
}

\setlist[itemize]{itemsep=0.3em,topsep=0.4em}
\setlist[enumerate]{itemsep=0.3em,topsep=0.4em}
\renewcommand{\arraystretch}{1.2}

\pagestyle{fancy}
\setlength{\headheight}{14pt}
\fancyhf{}
\lhead{StatProg II}
\rhead{Practical 3: Debugging}
\cfoot{\thepage}

\title{\textbf{Statistical Programming II}\\
Practical 3: Debugging\\
\large Exercise 1--5 答案与复习笔记}
\author{}
\date{}

\begin{document}

\maketitle
\tableofcontents
\newpage

\part{中文答案}

\section{Exercise 1：调试带列参数的函数}

\subsection{问题代码}

\begin{lstlisting}[style=rcode]
mean_mass_by <- function(data, group_var) {
  data |>
    group_by(group_var) |>
    summarise(
      mean_mass = mean(body_mass_g, na.rm = TRUE)
    )
}
\end{lstlisting}

如果调用：

\begin{lstlisting}[style=rcode]
mean_mass_by(penguins, species)
\end{lstlisting}

函数中的 \lstinline|group_by(group_var)| 不会自动把
\lstinline|group_var| 替换为调用时传入的列 \lstinline|species|。
\texttt{dplyr} 的 \lstinline|group_by()| 使用 data masking（数据掩码）和
tidy evaluation（整洁求值）。在自定义函数中转交一个未加引号的列名时，
需要明确告诉 \texttt{dplyr}：这里传入的是一列表达式。

\subsection{修正答案}

\begin{lstlisting}[style=rcode]
mean_mass_by <- function(data, group_var) {
  data |>
    group_by({{ group_var }}) |>
    summarise(
      mean_mass = mean(body_mass_g, na.rm = TRUE),
      .groups = "drop"
    )
}

mean_mass_by(penguins, species)
\end{lstlisting}

\lstinline|{{ group_var }}| 称为 embracing。可以把它理解为：

\begin{quote}
取得调用者传入的列名表达式，并把它放到 \lstinline|group_by()| 中求值。
\end{quote}

这里传参时不需要加引号：

\begin{lstlisting}[style=rcode]
mean_mass_by(penguins, species)
\end{lstlisting}

如果函数设计为接收字符串，例如 \lstinline|"species"|，则应写成：

\begin{lstlisting}[style=rcode]
mean_mass_by_string <- function(data, group_var) {
  data |>
    group_by(.data[[group_var]]) |>
    summarise(
      mean_mass = mean(body_mass_g, na.rm = TRUE),
      .groups = "drop"
    )
}

mean_mass_by_string(penguins, "species")
\end{lstlisting}

\subsection{本题要点}

\begin{itemize}
  \item 未加引号的列名参数：使用 \lstinline|{{ column }}|。
  \item 字符串形式的列名：使用 \lstinline|.data[[column]]|。
  \item \lstinline|na.rm = TRUE| 用于在计算均值时忽略缺失值。
  \item \lstinline|.groups = "drop"| 让结果不再保留分组结构。
\end{itemize}

\section{Exercise 2：错误发生在辅助函数内部}

\subsection{问题代码}

\begin{lstlisting}[style=rcode]
mean_body_mass <- function(x) {
  mean(x, na.rm = TREU)
}

summarise_species <- function(data) {
  data |>
    group_by(species) |>
    summarise(
      mean_mass = mean_body_mass(body_mass_g)
    )
}

summarise_species(penguins)
\end{lstlisting}

运行后会出现类似错误：

\begin{lstlisting}[style=textcode]
Error in `summarise()`:
...
Caused by error in `mean()`:
! object 'TREU' not found
\end{lstlisting}

\subsection{错误原因}

\lstinline|TREU| 是拼写错误。R 的逻辑值必须写成：

\begin{lstlisting}[style=rcode]
TRUE
FALSE
\end{lstlisting}

R 会把没有引号的 \lstinline|TREU| 当作对象名查找；当前环境中不存在这个
对象，因此产生 \lstinline|object 'TREU' not found|。

外层报错位置可能显示在 \lstinline|summarise()| 中，但真正出错的调用链是：

\begin{lstlisting}[style=textcode]
summarise_species()
  -> summarise()
     -> mean_body_mass()
        -> mean(..., na.rm = TREU)
\end{lstlisting}

因此，阅读嵌套错误时不能只看最外层函数，还应继续阅读
\texttt{Caused by error in ...} 部分。

\subsection{修正答案}

\begin{lstlisting}[style=rcode]
mean_body_mass <- function(x) {
  mean(x, na.rm = TRUE)
}

summarise_species <- function(data) {
  data |>
    group_by(species) |>
    summarise(
      mean_mass = mean_body_mass(body_mass_g),
      .groups = "drop"
    )
}

summarise_species(penguins)
\end{lstlisting}

\subsection{本题要点}

\begin{itemize}
  \item 报错显示在外层函数，不等于错误一定发生在外层函数。
  \item 应沿着 error message 中的调用关系寻找最内层原因。
  \item \lstinline|TRUE| 和 \lstinline|FALSE| 是 R 的保留逻辑值，不能随意拼写。
  \item 将重复计算封装为辅助函数没有问题，但调试时也必须检查辅助函数。
\end{itemize}

\section{Exercise 3：使用 \texttt{browser()} 调试递归函数}

\subsection{原始函数及递归逻辑}

\begin{lstlisting}[style=rcode]
my_factorial <- function(n) {
  if (n == 1) return(1)
  return(n * my_factorial(n - 1))
}
\end{lstlisting}

对于 \lstinline|my_factorial(5)|，递归过程是：

\begin{lstlisting}[style=textcode]
my_factorial(5)
= 5 * my_factorial(4)
= 5 * 4 * my_factorial(3)
= 5 * 4 * 3 * my_factorial(2)
= 5 * 4 * 3 * 2 * my_factorial(1)
= 5 * 4 * 3 * 2 * 1
= 120
\end{lstlisting}

\lstinline|if (n == 1) return(1)| 是 base case（终止条件）。如果递归过程
永远无法到达这个条件，函数就无法正常结束。

\subsection{加入 \texttt{browser()}}

\begin{lstlisting}[style=rcode]
my_factorial <- function(n) {
  browser()

  if (n == 1) return(1)
  return(n * my_factorial(n - 1))
}

my_factorial(0)
\end{lstlisting}

函数运行到 \lstinline|browser()| 时会暂停，并显示：

\begin{lstlisting}[style=textcode]
Browse[1]>
\end{lstlisting}

此时函数并没有结束。R 只是把执行暂停，让我们检查当前环境、逐行运行代码，
或者查看函数调用链。

\subsection{为什么不能只输入 \texttt{n}}

函数参数叫 \lstinline|n|，但在 debugger 中，单独的 \lstinline|n| 也是
\texttt{next} 命令。因此：

\begin{lstlisting}[style=rcode]
print(n)  # inspect the current value of n
n         # execute the debugger's "next" command
\end{lstlisting}

这解释了为什么输入 \lstinline|n| 后没有打印变量，反而看到了下一行代码。

\subsection{\texttt{my\_factorial(0)} 为什么无法结束}

第一次暂停时：

\begin{lstlisting}[style=textcode]
Browse[1]> print(n)
[1] 0
\end{lstlisting}

因为 \lstinline|0 == 1| 为 \lstinline|FALSE|，函数继续调用：

\begin{lstlisting}[style=rcode]
my_factorial(n - 1)
# my_factorial(-1)
\end{lstlisting}

之后参数按照下面的顺序变化：

\begin{lstlisting}[style=textcode]
0 -> -1 -> -2 -> -3 -> -4 -> ...
\end{lstlisting}

它永远不会等于 1，所以递归无法停止，最终通常会产生栈空间相关错误。
问题不在于数学上没有 \(0!\)。事实上：

\[
0! = 1.
\]

真正的问题是原函数没有把 0 纳入终止条件。

\subsection{\texttt{my\_factorial(3.5)} 为什么也不行}

\begin{lstlisting}[style=textcode]
3.5 -> 2.5 -> 1.5 -> 0.5 -> -0.5 -> ...
\end{lstlisting}

这一序列同样不会准确到达 1。此外，本题中的阶乘函数只应接受非负整数，
所以小数应在递归开始前直接被拒绝。

\subsection{常用调试命令}

\begin{longtable}{>{\ttfamily}p{0.16\textwidth}p{0.73\textwidth}}
\toprule
\normalfont 命令 & 含义 \\
\midrule
\endhead
print(n) & 查看变量 \texttt{n} 的当前值。\\
n & 执行下一条语句；通常不主动进入普通函数内部。\\
s & step into：执行下一步，并尝试进入该行调用的函数。\\
f & finish：尝试执行完当前函数，然后返回上一层。\\
c & continue：继续运行到下一个断点、错误或程序结束。\\
where & 显示当前调用栈。\\
Q & 立即退出当前调试过程。\\
\bottomrule
\end{longtable}

\subsection{如何理解 \texttt{where}}

例如：

\begin{lstlisting}[style=textcode]
where 1 at #4: my_factorial(n - 1)
where 2 at #4: my_factorial(n - 1)
where 3 at #4: my_factorial(n - 1)
where 4 at #4: my_factorial(n - 1)
where 5: my_factorial(0)
\end{lstlisting}

可以简化为：

\begin{lstlisting}[style=textcode]
my_factorial(0)
  -> my_factorial(-1)
     -> my_factorial(-2)
        -> my_factorial(-3)
           -> my_factorial(-4)  # current frame
\end{lstlisting}

\texttt{where 1} 是最靠近当前执行位置的调用，最后的
\texttt{where 5} 是最早的调用。外层函数都在等待更内层函数返回结果。

\texttt{Browse[1]} 中的数字 1 不是变量 \lstinline|n| 的值，也不是当前
递归深度。递归层数应通过 \lstinline|where| 判断。

\subsection{正确修复}

\begin{lstlisting}[style=rcode]
my_factorial <- function(n) {
  stopifnot(
    n >= 0,
    n == as.integer(n)
  )

  if (n <= 1) return(1)

  return(n * my_factorial(n - 1))
}

my_factorial(0)    # 1
my_factorial(1)    # 1
my_factorial(5)    # 120
my_factorial(3.5)  # Error
\end{lstlisting}

修复包含两部分：

\begin{enumerate}
  \item \lstinline|stopifnot()| 检查 \lstinline|n| 是否为非负整数；
  \item \lstinline|n <= 1| 同时正确处理 \(0!\) 和 \(1!\)。
\end{enumerate}

输入检查必须放在终止条件之前，否则负数也可能因为
\lstinline|n <= 1| 被错误地返回 1。调试完成后应删除函数中的
\lstinline|browser()|。

\section{Exercise 4：使用 \texttt{debug()} 和 \texttt{debugonce()}}

\subsection{问题函数}

\begin{lstlisting}[style=rcode]
standardise <- function(x) {
  (x - mean(x)) / sd(x)
}

standardise(c(1, 2, NA, 4, 5))
# NA NA NA NA NA
\end{lstlisting}

标准化公式为：

\[
z_i = \frac{x_i-\bar{x}}{s}.
\]

只有一个元素为 \lstinline|NA|，但结果全部变成了 \lstinline|NA|。

\subsection{使用 \texttt{debugonce()}}

\begin{lstlisting}[style=rcode]
debugonce(standardise)
standardise(c(1, 2, NA, 4, 5))
\end{lstlisting}

\lstinline|debugonce(standardise)| 不会立即运行函数。它表示：

\begin{quote}
下一次调用 \lstinline|standardise()| 时进入调试器；完成这一次调用后，
自动取消调试标记。
\end{quote}

进入 \lstinline|Browse[1]>| 后检查：

\begin{lstlisting}[style=rcode]
print(x)
# [1] 1 2 NA 4 5

mean(x)
# [1] NA

sd(x)
# [1] NA
\end{lstlisting}

因此原计算实际上变成：

\begin{lstlisting}[style=rcode]
(x - NA) / NA
\end{lstlisting}

缺失值参与算术运算时会继续传播，所以所有结果都成为
\lstinline|NA|。这称为 NA propagation。

\subsection{修正答案}

\begin{lstlisting}[style=rcode]
standardise <- function(x) {
  (x - mean(x, na.rm = TRUE)) /
    sd(x, na.rm = TRUE)
}

standardise(c(1, 2, NA, 4, 5))
# -1.095445 -0.547723 NA 0.547723 1.095445
\end{lstlisting}

计算均值和标准差时使用的是：

\begin{lstlisting}[style=textcode]
1, 2, 4, 5
\end{lstlisting}

它们的均值为 3。原本缺失的第三个元素仍然是 \lstinline|NA|。
\lstinline|na.rm = TRUE| 只是让 \lstinline|mean()| 和 \lstinline|sd()|
在计算时忽略缺失值，并不会自动填补原数据。

\subsection{\texttt{debug()} 与 \texttt{debugonce()} 的区别}

\begin{lstlisting}[style=rcode]
debug(standardise)
\end{lstlisting}

使用 \lstinline|debug()| 后，每次调用函数都会进入调试器，直到手动取消：

\begin{lstlisting}[style=rcode]
undebug(standardise)
\end{lstlisting}

也可以检查函数当前是否带有调试标记：

\begin{lstlisting}[style=rcode]
isdebugged(standardise)
# TRUE or FALSE
\end{lstlisting}

\begin{longtable}{p{0.24\textwidth}p{0.25\textwidth}p{0.20\textwidth}p{0.20\textwidth}}
\toprule
方法 & 设置位置 & 生效范围 & 取消方法 \\
\midrule
\endhead
\lstinline|browser()| & 写在函数内部 & 每次执行到该行 & 从函数中删除 \lstinline|browser()| \\
\lstinline|debugonce(f)| & 函数外部 & 只调试下一次调用 & 自动取消 \\
\lstinline|debug(f)| & 函数外部 & 调试之后的每次调用 & \lstinline|undebug(f)| \\
\bottomrule
\end{longtable}

\section{Exercise 5：调试 Quarto 文档}

\subsection{调试思路}

Quarto 文档从上到下执行代码块。某个 chunk 失败以后，后面的 chunk
通常不会继续执行。因此应采用以下过程：

\begin{enumerate}
  \item 找到报错信息中第一个真正的错误；
  \item 确认出错文件、行号和 chunk；
  \item 修复当前错误；
  \item 重新渲染，让下一个错误显现；
  \item 重复以上步骤，直到渲染成功。
\end{enumerate}

\subsection{错误一：用 \texttt{=} 代替 \texttt{==}}

错误代码：

\begin{lstlisting}[style=rcode]
penguins |>
  filter(species = "Adelie")
\end{lstlisting}

\lstinline|=| 在函数调用中用于给参数命名，不表示比较。筛选
\lstinline|species| 等于 \lstinline|"Adelie"| 的行应写成：

\begin{lstlisting}[style=rcode]
penguins |>
  filter(species == "Adelie")
\end{lstlisting}

\begin{lstlisting}[style=textcode]
species = "Adelie"   # assignment/naming
species == "Adelie"  # equality comparison
\end{lstlisting}

这是 R/\texttt{dplyr} 代码执行阶段的错误。

\subsection{错误二：没有加载 \texttt{ggplot2}}

下列函数来自 \texttt{ggplot2}：

\begin{lstlisting}[style=rcode]
ggplot()
aes()
geom_point()
\end{lstlisting}

如果文档中没有加载该包，干净的渲染 session 会出现类似：

\begin{lstlisting}[style=textcode]
could not find function "ggplot"
\end{lstlisting}

因此应在文档中明确写出：

\begin{lstlisting}[style=rcode]
library(palmerpenguins)
library(dplyr)
library(ggplot2)
\end{lstlisting}

代码在 Console 中能运行，不代表它在 Quarto 中一定能运行。Console
可能已经加载过 \texttt{tidyverse}，而 Quarto 渲染通常在新的 R session
中从头执行。文档必须自己加载所需的包，不能依赖交互式环境中的历史状态。

\subsection{错误三：\texttt{echo=TREU}}

错误的 chunk header：

\begin{lstlisting}[style=textcode]
```{r, echo=TREU}
summary(penguins)
```
\end{lstlisting}

\lstinline|TREU| 不是 R 的逻辑值。可以改为：

\begin{lstlisting}[style=textcode]
```{r, echo=TRUE}
summary(penguins)
```
\end{lstlisting}

更推荐使用 Quarto 的 YAML 风格 chunk option：

\begin{lstlisting}[style=textcode]
```{r}
#| echo: true

summary(penguins)
```
\end{lstlisting}

\lstinline|echo: true| 表示运行代码并在最终文档中显示源代码；
\lstinline|echo: false| 表示运行代码但隐藏源代码。它本身不决定是否运行
这个 chunk。

该错误发生在 knitr 解析 chunk option 的阶段，甚至可能早于
\lstinline|summary(penguins)| 的实际执行。

\subsection{完整修复后的 \texttt{broken.qmd}}

\begin{lstlisting}[style=textcode]
---
title: "Penguin bill measurements"
format: html
---

## First chunk

```{r}
library(palmerpenguins)
library(dplyr)
library(ggplot2)

head(penguins)
```

## Second chunk

```{r}
penguins |>
  filter(species == "Adelie") |>
  ggplot(aes(
    x = bill_length_mm,
    y = bill_depth_mm
  )) +
  geom_point()
```

## Third chunk

```{r}
#| echo: true

summary(penguins)
```
\end{lstlisting}

在文件所在目录运行：

\begin{lstlisting}[style=textcode]
quarto render broken.qmd
\end{lstlisting}

成功后会生成 HTML 文件。

\subsection{判断错误属于哪一层}

\begin{longtable}{p{0.48\textwidth}p{0.40\textwidth}}
\toprule
报错线索 & 可能的问题层 \\
\midrule
\endhead
\texttt{Error in filter()}、\texttt{object not found}、\texttt{could not find function} & R 代码 \\
\texttt{Quitting from ... [chunk-name]} & knitr 执行 chunk 时停止；根本原因可能仍是 R 错误 \\
错误的 \texttt{echo}、\texttt{eval} 等 chunk option & knitr \\
\texttt{YAMLException} 或 front matter 格式错误 & YAML/Quarto \\
Pandoc conversion 或 extension 相关错误 & Quarto/Pandoc \\
\texttt{LaTeX Error}、缺少 \texttt{.sty} 文件 & LaTeX/PDF 编译 \\
\bottomrule
\end{longtable}

给 chunk 添加名称可以提高可定位性：

\begin{lstlisting}[style=textcode]
```{r}
#| label: plot-adelie

...
```
\end{lstlisting}

此时报错会显示 \lstinline|[plot-adelie]|，比
\lstinline|[unnamed-chunk-2]| 更容易判断位置。

\subsection{HTML 与 PDF 渲染的差别}

HTML 渲染大致经过：

\begin{lstlisting}[style=textcode]
.qmd -> knitr runs R -> Pandoc -> HTML
\end{lstlisting}

PDF 渲染还需要 LaTeX：

\begin{lstlisting}[style=textcode]
.qmd -> knitr runs R -> Pandoc -> LaTeX -> PDF
\end{lstlisting}

因此，把 \lstinline|format: html| 改为 \lstinline|format: pdf| 后出现
\lstinline|No TeX installation was detected|，并不说明 R 代码有错，而是
PDF 编译环境缺失。本题的核心是：

\begin{quote}
找到第一个真正的错误，判断它来自哪一层，修复后重新渲染；不要只盯着最后
一行 \texttt{Execution halted}。
\end{quote}

\newpage
\part{Deutsche Fassung}

\section{Exercise 1: Spaltenargumente mit tidy evaluation}

\subsection{Problem}

\begin{lstlisting}[style=rcode]
mean_mass_by <- function(data, group_var) {
  data |>
    group_by(group_var) |>
    summarise(
      mean_mass = mean(body_mass_g, na.rm = TRUE)
    )
}
\end{lstlisting}

\lstinline|group_by(group_var)| verwendet das Argument nicht automatisch als
die beim Funktionsaufruf übergebene Spalte. \texttt{dplyr} arbeitet hier mit
data masking und tidy evaluation.

\subsection{Korrektur}

\begin{lstlisting}[style=rcode]
mean_mass_by <- function(data, group_var) {
  data |>
    group_by({{ group_var }}) |>
    summarise(
      mean_mass = mean(body_mass_g, na.rm = TRUE),
      .groups = "drop"
    )
}

mean_mass_by(penguins, species)
\end{lstlisting}

\lstinline|{{ group_var }}| übernimmt den nicht als String übergebenen
Spaltenausdruck aus dem Funktionsaufruf. Für einen Spaltennamen als String
würde man stattdessen \lstinline|.data[[group_var]]| verwenden.

\section{Exercise 2: Fehler in einer Hilfsfunktion}

\subsection{Fehlerursache}

\begin{lstlisting}[style=rcode]
mean_body_mass <- function(x) {
  mean(x, na.rm = TREU)
}
\end{lstlisting}

\lstinline|TREU| ist kein logischer Wert in R. R sucht deshalb nach einem
Objekt mit diesem Namen und meldet:

\begin{lstlisting}[style=textcode]
object 'TREU' not found
\end{lstlisting}

Der Fehler wird möglicherweise zunächst von \lstinline|summarise()| gemeldet,
entsteht aber tatsächlich im inneren Aufruf von \lstinline|mean()|.

\subsection{Korrektur}

\begin{lstlisting}[style=rcode]
mean_body_mass <- function(x) {
  mean(x, na.rm = TRUE)
}

summarise_species <- function(data) {
  data |>
    group_by(species) |>
    summarise(
      mean_mass = mean_body_mass(body_mass_g),
      .groups = "drop"
    )
}
\end{lstlisting}

Bei verschachtelten Fehlermeldungen sollte deshalb insbesondere der Abschnitt
\texttt{Caused by error in ...} gelesen werden.

\section{Exercise 3: Rekursion mit \texttt{browser()} untersuchen}

\subsection{Warum die Funktion fehlschlägt}

\begin{lstlisting}[style=rcode]
my_factorial <- function(n) {
  if (n == 1) return(1)
  return(n * my_factorial(n - 1))
}
\end{lstlisting}

Für \lstinline|my_factorial(5)| erreicht die Rekursion den Abbruchfall
\lstinline|n == 1|. Für \lstinline|my_factorial(0)| entsteht dagegen:

\begin{lstlisting}[style=textcode]
0 -> -1 -> -2 -> -3 -> ...
\end{lstlisting}

Für \lstinline|my_factorial(3.5)| gilt entsprechend:

\begin{lstlisting}[style=textcode]
3.5 -> 2.5 -> 1.5 -> 0.5 -> -0.5 -> ...
\end{lstlisting}

In beiden Fällen wird 1 niemals erreicht.

\subsection{Interaktives Debugging}

\begin{lstlisting}[style=rcode]
my_factorial <- function(n) {
  browser()
  if (n == 1) return(1)
  return(n * my_factorial(n - 1))
}
\end{lstlisting}

\texttt{Browse[1]>} bedeutet, dass R die Ausführung angehalten hat. Wichtig ist
der Namenskonflikt:

\begin{lstlisting}[style=rcode]
print(n)  # Wert der Variablen anzeigen
n         # Debugger-Befehl "next"
\end{lstlisting}

\begin{longtable}{>{\ttfamily}p{0.16\textwidth}p{0.73\textwidth}}
\toprule
\normalfont Befehl & Bedeutung \\
\midrule
\endhead
n & Nächste Anweisung ausführen.\\
s & In einen Funktionsaufruf hineingehen.\\
f & Aktuelle Funktion möglichst bis zum Ende ausführen.\\
c & Bis zum nächsten Haltepunkt oder Fehler weiterlaufen.\\
where & Den aktuellen Aufrufstapel anzeigen.\\
Q & Das Debugging sofort abbrechen.\\
\bottomrule
\end{longtable}

Mehrere Einträge bei \lstinline|where| zeigen mehrere noch nicht abgeschlossene
rekursive Aufrufe. Die äußeren Aufrufe warten jeweils auf das Ergebnis des
inneren Aufrufs. Die Zahl in \texttt{Browse[1]} ist weder der Wert von
\lstinline|n| noch die Rekursionstiefe.

\subsection{Korrektur}

\begin{lstlisting}[style=rcode]
my_factorial <- function(n) {
  stopifnot(
    n >= 0,
    n == as.integer(n)
  )

  if (n <= 1) return(1)

  return(n * my_factorial(n - 1))
}
\end{lstlisting}

\lstinline|stopifnot()| lehnt negative Zahlen und Dezimalzahlen ab.
\lstinline|n <= 1| behandelt sowohl \(0!\) als auch \(1!\) korrekt. Nach dem
Debugging muss \lstinline|browser()| entfernt werden.

\section{Exercise 4: \texttt{debug()} und \texttt{debugonce()}}

\subsection{Fehleranalyse}

\begin{lstlisting}[style=rcode]
standardise <- function(x) {
  (x - mean(x)) / sd(x)
}

debugonce(standardise)
standardise(c(1, 2, NA, 4, 5))
\end{lstlisting}

Im Debugger ergibt sich:

\begin{lstlisting}[style=rcode]
mean(x)  # NA
sd(x)    # NA
\end{lstlisting}

Damit wird faktisch \lstinline|(x - NA) / NA| berechnet. Durch die
Weitergabe fehlender Werte werden alle Ergebnisse zu \lstinline|NA|.

\subsection{Korrektur}

\begin{lstlisting}[style=rcode]
standardise <- function(x) {
  (x - mean(x, na.rm = TRUE)) /
    sd(x, na.rm = TRUE)
}
\end{lstlisting}

Der ursprünglich fehlende Wert bleibt \lstinline|NA|. Er wird lediglich bei
der Berechnung von Mittelwert und Standardabweichung ignoriert.

\begin{longtable}{p{0.24\textwidth}p{0.27\textwidth}p{0.28\textwidth}}
\toprule
Werkzeug & Verhalten & Beenden \\
\midrule
\endhead
\lstinline|browser()| & Haltepunkt direkt im Funktionskörper & Zeile aus der Funktion entfernen \\
\lstinline|debugonce(f)| & Nur der nächste Aufruf wird debuggt & automatische Aufhebung \\
\lstinline|debug(f)| & Jeder weitere Aufruf wird debuggt & \lstinline|undebug(f)| \\
\bottomrule
\end{longtable}

\section{Exercise 5: Ein Quarto-Dokument debuggen}

\subsection{Die drei Fehler}

\begin{enumerate}
  \item In \lstinline|filter(species = "Adelie")| wurde eine Argumentzuweisung
  statt eines Vergleichs verwendet. Richtig ist:

\begin{lstlisting}[style=rcode]
filter(species == "Adelie")
\end{lstlisting}

  \item \lstinline|ggplot()|, \lstinline|aes()| und
  \lstinline|geom_point()| stammen aus \texttt{ggplot2}. Das Paket muss im
  Dokument explizit geladen werden:

\begin{lstlisting}[style=rcode]
library(ggplot2)
\end{lstlisting}

  \item \lstinline|echo=TREU| enthält erneut den Tippfehler
  \lstinline|TREU|. Empfohlen ist die Quarto-Schreibweise:

\begin{lstlisting}[style=textcode]
```{r}
#| echo: true
summary(penguins)
```
\end{lstlisting}
\end{enumerate}

\subsection{Korrigiertes Dokument}

\begin{lstlisting}[style=textcode]
---
title: "Penguin bill measurements"
format: html
---

```{r}
library(palmerpenguins)
library(dplyr)
library(ggplot2)
```

```{r}
penguins |>
  filter(species == "Adelie") |>
  ggplot(aes(
    x = bill_length_mm,
    y = bill_depth_mm
  )) +
  geom_point()
```

```{r}
#| echo: true
summary(penguins)
```
\end{lstlisting}

Gerendert wird mit:

\begin{lstlisting}[style=textcode]
quarto render broken.qmd
\end{lstlisting}

\subsection{Fehlerebenen}

\begin{longtable}{p{0.48\textwidth}p{0.40\textwidth}}
\toprule
Hinweis & Wahrscheinliche Ebene \\
\midrule
\endhead
\texttt{Error in filter()}, \texttt{object not found} & R \\
\texttt{Quitting from ... [chunk-name]} & knitr stoppt beim Chunk \\
Fehlerhafte Chunk-Option & knitr \\
YAML-Fehler & YAML/Quarto \\
Konvertierungsfehler & Quarto/Pandoc \\
\texttt{LaTeX Error}, fehlende \texttt{.sty}-Datei & LaTeX \\
\bottomrule
\end{longtable}

Quarto führt Chunks von oben nach unten aus. Deshalb wird zunächst nur der erste
Fehler sichtbar. Dieser wird korrigiert, anschließend wird erneut gerendert.
Ein PDF-Render benötigt zusätzlich eine LaTeX-Installation; ein fehlendes
TeX-System ist kein R-Fehler.

\section*{Kurzfazit}
\addcontentsline{toc}{section}{Kurzfazit}

\begin{itemize}
  \item Bei \texttt{dplyr}-Funktionen muss die Auswertung von
  Spaltenargumenten bewusst behandelt werden.
  \item Verschachtelte Fehlermeldungen müssen bis zur inneren Ursache gelesen
  werden.
  \item \texttt{browser()}, \texttt{debugonce()} und \texttt{debug()} halten
  Funktionen auf unterschiedliche Weise an, verwenden aber dieselben
  Debugger-Befehle.
  \item Rekursive Funktionen benötigen einen erreichbaren Abbruchfall und eine
  Prüfung ihres Eingabebereichs.
  \item Beim Quarto-Rendering muss zwischen R, knitr, Quarto/Pandoc und LaTeX
  unterschieden werden.
\end{itemize}

\end{document}
